\documentclass[11pt]{amsart}

\usepackage[T1]{fontenc}
\usepackage{lmodern}
\usepackage{microtype}
\usepackage{mathtools}
\usepackage{amssymb}
\usepackage{booktabs}
\usepackage{array}
\usepackage{enumitem}
\usepackage{xcolor}
\usepackage{hyperref}
\usepackage[nameinlink,noabbrev]{cleveref}
\usepackage{listings}
\usepackage{geometry}
\geometry{margin=1.08in}

\hypersetup{
  colorlinks=true,
  linkcolor=blue!45!black,
  citecolor=green!35!black,
  urlcolor=blue!55!black,
  pdfauthor={Mansehej Singh},
  pdftitle={Verification supplement: Exact non-scattering homogeneous dielectrics from Schiffer counterexamples}
}

\newtheorem{theorem}{Theorem}[section]
\newtheorem{proposition}[theorem]{Proposition}
\newtheorem{lemma}[theorem]{Lemma}
\newtheorem{remark}[theorem]{Remark}

\newcommand{\D}{\mathbb D}
\newcommand{\R}{\mathbb R}
\newcommand{\norm}[1]{\left\lVert #1\right\rVert}
\newcommand{\abs}[1]{\left\lvert #1\right\rvert}
\newcommand{\Rea}{\operatorname{Re}}

\title[Verification supplement]
{Verification supplement for\\
\emph{Exact non-scattering homogeneous dielectrics from Schiffer counterexamples}}
\author{Mansehej Singh}
\date{19 August 2026}

\begin{document}

\begin{abstract}
This supplement records the computer-assisted and formal-verification layer
supporting the accompanying paper.  It specifies the imported
Colbrook--Stepaniants theorem, derives the conservative preconditioner bound
$\norm{\mathcal A}<1180$, gives the detailed Bessel perturbation estimates,
normalizes the endpoint inequalities by exact rational arithmetic, and
documents the independent executable verifiers and the narrow Lean 4
re-verification of the certificate arithmetic.
The supplement is deliberately separated from the research article: the
article presents the mathematical discovery and proof, while this document
provides the audit trail needed to replay and challenge the certificate.
\end{abstract}

\maketitle

\section{Logical structure and scope}

The proof has three layers.
\begin{enumerate}[label=\textup{(\roman*)},leftmargin=2.2em]
  \item The theorem of Colbrook and Stepaniants
  \cite{ColbrookStepaniants2026} supplies an exact noncircular Schiffer zero
  in a weighted coefficient algebra, together with an invertible
  preconditioner, contraction bounds, infinite-tail control, and uniform
  geometry estimates.
  \item The new analytic layer defines the Bessel-deformed equation
  $F_\tau=0$, derives the value and derivative perturbation estimates, and
  proves that the same contraction ball remains valid for
  $0\le\tau\le10^{-13}$.
  \item The physical layer pushes the exact fixed point through the conformal
  map and a dilation, yielding a constant-index transmission solution whose
  exterior scattered field is identically zero.
\end{enumerate}

The large upstream interval computation is not reimplemented in this
supplement.  Instead, the released upstream certificate is treated as a
pinned theorem dependency and replayed from its frozen source and inverse.
The new scalar inequalities are checked independently by exact-rational
Python, Node/BigInt arithmetic, and Lean normalization.

\section{Pinned upstream theorem}

The upstream fixed-disc equation is
\begin{equation}\label{eq:F0}
  F_0(\gamma,p)=\gamma+\abs p^2(1+K\gamma)=0,
\end{equation}
with ten-fold shape series
\[
  p(z)=\sum_{j\ge0}p_jz^{10j},
  \qquad
  \phi_p(z)=z+\sum_{j\ge1}
  \frac{p_j}{p_0(10j+1)}z^{10j+1}.
\]
The certificate parameters are
\[
  (L,S,R,J,\rho)=(60,40,30,30,1.05),
  \qquad N=2471,
\]
and the validated radius is
\[
  r=10^{-6}.
\]
The authenticated coarse majorants are
\begin{equation}\label{eq:coarse}
  Y\le1.59\times10^{-10},
  \qquad Z\le0.621,
  \qquad C_2\le122,
  \qquad C_3\le0.012.
\end{equation}
They imply
\begin{align}
  R_0&:=Y+Zr+C_2r^2+C_3r^3-r
  \le-0.000000378718999999988,
  \label{eq:R0}\\
  L_0&:=Z+2C_2r+3C_3r^2
  \le0.621244000000036.
  \label{eq:L0}
\end{align}
The same theorem proves on the whole validated ball
\begin{equation}\label{eq:geometry}
  \norm p_\rho<55,
  \qquad 31<p_0<32,
  \qquad
  \sum_{j\ge1}\abs{p_j/p_0}<0.65,
\end{equation}
and keeps the first nonconstant shape coefficient away from zero.  Hence
$\Rea\phi_p'>0.35$ on the closed unit disc, so $\phi_p$ is univalent on a
complex collar and its image is a noncircular analytic Jordan domain.

The release pins the following upstream objects:
\begin{itemize}[leftmargin=2em]
  \item repository and commit;
  \item exact hexadecimal center;
  \item frozen binary64 approximate inverse;
  \item source manifest and certificate schema;
  \item authenticated 192-bit and 256-bit MPFR audit hashes.
\end{itemize}
The concrete values are stored in
\path{verification/upstream.lock} and
\path{verification/upstream_snapshot.json}.

\section{Hardening the preconditioner bound}

Let
\[
  b=\bar p_0,
  \qquad
  \alpha=\max(1,\norm{\mathcal A}).
\]
The upstream nonlinear estimate defines
\begin{equation}\label{eq:C3}
  C_3=\frac{\alpha}{96b^2}.
\end{equation}
The exact first shape token is
\[
  b=\texttt{0x1.ff78dc9f663aap+4}
  =\frac{4498956314423765}{140737488355328}<32.
\]
Combining this exact inequality with the simple authenticated majorant
$C_3\le0.012$ yields
\begin{align}
  \alpha
  &=96b^2C_3\notag\\
  &<96\cdot32^2\cdot0.012
   =1179.648<1180.
  \label{eq:alpha}
\end{align}
Thus the extension proof needs neither the full inverse norm nor a long
decimal reconstructed from the explanatory paper.  The executable
certificate uses the integer upper bound
\begin{equation}\label{eq:alpha1180}
  \boxed{\norm{\mathcal A}<1180.}
\end{equation}

\section{Detailed Bessel perturbation estimates}

For $\tau\ge0$, set
\begin{equation}\label{eq:H}
  H_{\tau,p}
  =\sum_{n=0}^{\infty}
  \frac{(-1)^n}{(n!)^2}
  \left(\frac{\tau\abs{\phi_p}^2}{4}\right)^n
\end{equation}
and
\begin{equation}\label{eq:Ftau}
  F_\tau(\gamma,p)
  =\gamma+\abs p^2K\gamma
  +\abs p^2\left(1-\frac\tau{p_0^2}\right)H_{\tau,p}.
\end{equation}
Write $\delta F_\tau=F_\tau-F_0$.

\subsection{Bound on the conformal radius factor}

Write $\phi_p(z)=z\psi_p(z)$.  From the coefficient formula and
\eqref{eq:geometry},
\begin{align}
  \norm{\psi_p}
  &\le1+\sum_{j\ge1}
  \frac{\abs{p_j}}{p_0(10j+1)}\rho^j\notag\\
  &\le1+\frac{\norm p_\rho-p_0}{11p_0}
  <1+\frac{55-31}{11\cdot31}
  <\frac{11}{10}.
  \label{eq:psi}
\end{align}
The upstream algebra norm is a weighted $\ell^1$ sum over a basis indexed
by an angular block index $\ell$ and a radial index $s$, with weight
$c_\ell\rho^\ell$ depending only on $\ell$, and multiplication by the
radial factor $\abs z^2=r^2$ acts as the pure radial shift
$(\ell,s)\mapsto(\ell,s+1)$
\cite[Validation Specification, \S\S2 and~4]{ColbrookStepaniantsCert2026}.
Multiplication by $\abs z^2$ is therefore nonexpansive, and
\begin{equation}\label{eq:R}
  \norm{\abs{\phi_p}^2}
  =\norm{\abs z^2\abs{\psi_p}^2}
  \le\norm{\psi_p}^2
  <\frac{121}{100}<\frac54.
\end{equation}
This radial-shift property is the single algebra fact imported beyond
submultiplicativity; it is restated as a lemma in the main article and is
item one of the recommended independent audits.

\subsection{Value perturbation}

Put
\[
  u=\frac\tau4\norm{\abs{\phi_p}^2}.
\]
For $0\le\tau\le10^{-13}$, $u<1/2$.  Since
$1/(n!)^2\le1$, the full Bessel series satisfies
\begin{equation}\label{eq:Hnorm}
  \norm{H_{\tau,p}}
  \le\sum_{n\ge0}u^n
  =\frac1{1-u}\le2.
\end{equation}
Similarly,
\begin{equation}\label{eq:Hminus1}
  \norm{H_{\tau,p}-1}
  \le\frac{u}{1-u}
  \le2u
  \le\frac58\tau.
\end{equation}
Hence
\begin{align}
 \norm{\left(1-\frac\tau{p_0^2}\right)H_{\tau,p}-1}
 &\le\norm{H_{\tau,p}-1}
 +\frac\tau{p_0^2}\norm{H_{\tau,p}}\notag\\
 &\le\left(\frac58+\frac2{31^2}\right)\tau
 <\frac{63}{100}\tau.
 \label{eq:B}
\end{align}
Using $\norm p^2<55^2=3025$,
\begin{equation}\label{eq:value}
  \norm{\delta F_\tau}
  <3025\cdot\frac{63}{100}\tau
  =1905.75\tau
  <1906\tau.
\end{equation}

\subsection{Derivative perturbation}

Let $\xi$ be a shape variation with full $X_\rho$ norm at most one.  The
upstream product norm weights the shape block by $2\bar p_0\rho^j$
\cite[Validation Specification, Eq.~(13)]{ColbrookStepaniantsCert2026},
which gives
\begin{equation}\label{eq:xi}
  \norm\xi\le\frac1{2b}<\frac1{62},
  \qquad
  \abs{\xi_0}<\frac1{62}.
\end{equation}
Differentiating $\psi_p$ gives
\begin{align}
  \norm{D\psi_p[\xi]}
  &\le\frac1{11}
  \left(\frac{\norm\xi}{p_0}
  +\frac{\norm p\abs{\xi_0}}{p_0^2}\right)\notag\\
  &<\frac1{11}
  \left(\frac1{62\cdot31}
  +\frac{55}{62\cdot31^2}\right)
  <\frac1{7000}.
  \label{eq:Dpsi}
\end{align}
Thus
\begin{equation}\label{eq:DR}
  \norm{D\abs{\phi_p}^2[\xi]}
  \le2\norm{\psi_p}\norm{D\psi_p[\xi]}
  <\frac1{3000},
\end{equation}
using the radial-shift property once more for the factor $\abs z^2$.
Termwise differentiation of \eqref{eq:H} and
$\sum_{n\ge1}nu^{n-1}=(1-u)^{-2}\le4$ give
\begin{equation}\label{eq:DH}
  \norm{DH_{\tau,p}[\xi]}
  \le\frac\tau4\norm{D\abs{\phi_p}^2[\xi]}
  \sum_{n\ge1}nu^{n-1}
  <\frac\tau{3000}.
\end{equation}

Let
\[
  B_{\tau,p}
  =\left(1-\frac\tau{p_0^2}\right)H_{\tau,p}-1.
\]
Then
\begin{align*}
  D\delta F_\tau[\xi]
  &=2\Rea(\bar p\xi)B_{\tau,p}\\
  &\quad+\abs p^2\left[
    \left(1-\frac\tau{p_0^2}\right)DH_{\tau,p}[\xi]
    +\frac{2\tau\xi_0}{p_0^3}H_{\tau,p}
  \right].
\end{align*}
Applying \eqref{eq:B}, \eqref{eq:Hnorm}, \eqref{eq:xi}, and
\eqref{eq:DH},
\begin{align}
  \norm{D\delta F_\tau}
  &<\left[
  \frac{110}{62}\frac{63}{100}
  +\frac{3025}{3000}
  +\frac{12100}{62\cdot31^3}
  \right]\tau\notag\\
  &<3\tau.
  \label{eq:derivative}
\end{align}
The last strict inequality is checked exactly in all three arithmetic
implementations.

\section{Exact closure of the contraction ball}

Let
\[
  T_\tau=I-\mathcal AF_\tau.
\]
The value perturbation changes the invariant-ball radii polynomial by at
most
\[
  1180\cdot1906\,\tau,
\]
and the derivative perturbation changes the contraction constant by at most
\[
  1180\cdot3\,\tau.
\]
Thus
\begin{align}
  R_\tau&\le R_0+1180\cdot1906\,\tau,
  \label{eq:Rtau}\\
  L_\tau&\le L_0+1180\cdot3\,\tau.
  \label{eq:Ltau}
\end{align}
At $\tau_*=10^{-13}$, exact rational arithmetic yields
\begin{align}
  R_{\tau_*}
  &\le-\frac{153810999999988}{10^{21}}
  =-1.53810999999988\times10^{-7}<0,
  \label{eq:Rend}\\
  L_{\tau_*}
  &\le\frac{621244000354036}{10^{15}}
  =0.621244000354036<1.
  \label{eq:Lend}
\end{align}
Because the right-hand sides are monotone in $\tau$, the same conclusions
hold for the whole interval $0\le\tau\le10^{-13}$.

For comparison, the explicit example $\tau=10^{-14}$ satisfies
\begin{align*}
  R_\tau
  &\le-3.56228199999988\times10^{-7}<0,\\
  L_\tau
  &\le0.621244000035436<1.
\end{align*}

\section{Executable certificate}

The authoritative new input is
\path{verification/extension_certificate.json}.  It contains the
rational constants, the parameter interval, the exact hexadecimal center
coefficient used in \eqref{eq:alpha}, and hashes of the pinned upstream
snapshot.  The primary checker
\path{verification/verify_extension.py}:
\begin{enumerate}[leftmargin=2.2em]
  \item parses all decimal values as exact rational numbers;
  \item cross-checks the snapshot hashes, repository, and commit against
  the pinned \path{upstream.lock} (file-level hash verification is
  performed by the upstream replay script below);
  \item parses the hexadecimal binary64 center token exactly and derives
  $b<32$, the ball bound $b+r/(2\cdot31)<32$, and $\alpha<1180$ rather
  than trusting them as loose text constants;
  \item rederives the perturbation-constant chain from the certified ball
  bounds $\norm p_\rho<55$ and $p_0>31$: the bounds $11/10$ for
  $\psi_p$, $5/4$ for $\abs{\phi_p}^2$, $63/100$ for the datum
  deviation, $1/7000$ and $1/3000$ for the derivatives, and finally
  $55^2\cdot63/100<1906$ and the three-term derivative sum below $3$;
  \item verifies the Bessel-series condition at the uniform endpoint
  $\tau=10^{-13}$;
  \item recomputes the radii polynomial, the contraction constant, and
  both endpoint inequalities, comparing each against the stored exact
  decimal expansions;
  \item fails by explicit exceptions and nonzero exit status.
\end{enumerate}
Proof conditions do not use Python's removable \texttt{assert} statement, so
normal and optimized execution are byte-identical.

A separately implemented Node program,
\path{verification/verify_extension.mjs}, uses BigInt rational arithmetic,
reads the same authoritative certificate, and replays the complete check
set above, including the derivation chain and the lock cross-check.

\subsection{Upstream replay}

\path{verification/fetch_and_check_upstream.sh} downloads the archived
upstream release from Zenodo (record 21765287), verifies the archive
digest recorded in \path{upstream.lock}, and then runs
\path{verification/cross_check_upstream.py}, which
\begin{enumerate}[leftmargin=2.2em]
  \item verifies the pinned SHA-256 hashes of the upstream center, inverse,
  certificate, and source manifest;
  \item compares the upstream certificate values
  $(Y,Z,C_2,C_3,r,\rho,L,S,R,J)$, the univalence sum, and the first-shape-%
  coefficient bound against \path{upstream_snapshot.json};
  \item rederives the ball bounds $31<p_0<32$ and $\norm p_\rho<55$ from
  the exact hexadecimal center coefficients plus the validated radius.
\end{enumerate}
Finally it replays the upstream authenticated checker
\path{verify_certificate.py} itself.  The successful run of this entire
pipeline from a clean fetch is recorded in
\path{verification/upstream_check.log}; the upstream checker's exact
radii-polynomial output reproduces \eqref{eq:R0} digit for digit.

\subsection{Adversarial tests}

The regression suite deliberately corrupts, one at a time: the
preconditioner bound; the value-perturbation constant; the conformal-square
majorant $5/4$; the endpoint $\tau$; a stored exact decimal expansion; the
upstream $Y$, $Z$, and $C_3$ majorants; the certified shape-norm ball bound
$55$; the exact hexadecimal center token; a pinned hash (forcing a
snapshot/lock disagreement); and the presence of the snapshot and lock
files themselves.  Every corruption must be rejected by the Python checker
\emph{and}, when Node is available, by the independent BigInt checker.
The suite also requires byte-equal output under \texttt{python3} and
\texttt{python3 -O}.

\section{Lean re-verification of the certificate arithmetic}

The release contains a small Lean 4 project pinned to Lean/mathlib
v4.32.1.  Its scope is deliberately narrow and is stated here exactly.

\begin{center}
\begin{tabular}{@{}>{\raggedright\arraybackslash}p{0.28\linewidth}>{\raggedright\arraybackslash}p{0.64\linewidth}@{}}
\toprule
Module & Source content \\
\midrule
\texttt{Certificate} & Exact rational evaluation of
\eqref{eq:R0}--\eqref{eq:Lend}, the hardened
$C_3,b\Rightarrow\alpha<1180$ implication, uniform closure over the whole
parameter interval, and the two scalar majorant inequalities
$55^2\cdot63/100<1906$ and the three-term derivative sum below $3$.\\
\texttt{FixedPoint} & Banach fixed-point existence and uniqueness on a
complete metric space, and the stability estimate for fixed points of two
nearby contractions, stated generically.\\
\texttt{LinearRegularity} & If $A$ and $A\circ B$ are bijective, so is
$B$.\\
\texttt{Scaling} & The scalar contrast computation $q=p_0^2/\tau$, its
positivity, and $q>1$ on the certified interval.\\
\bottomrule
\end{tabular}
\end{center}

What the Lean development is \emph{not}: it does not formalize the
operator-level perturbation estimates behind the constants $1906$ and $3$,
the weighted coefficient algebra, the upstream MPFR certificate, or any
planar scattering analysis (the Bessel--Herglotz identity, conformal
covariance, trace transformation, exterior uniqueness).  The generic
fixed-point lemmas are stated in the form the paper uses but are not
instantiated on the concrete coefficient spaces.  The accurate public
description is therefore:
\begin{quote}
A Lean re-verification of the exact rational certificate arithmetic,
together with generic fixed-point and invertibility lemmas in the form used
by the paper.  It is a third independent checker of the scalar arithmetic,
not a formalization of the non-scattering theorem.
\end{quote}

A static source scan (\path{formalization/scripts/check_no_holes.py})
rejects \texttt{sorry}, \texttt{admit}, \texttt{axiom},
\texttt{native\_decide}, \texttt{opaque}, \texttt{unsafe},
\texttt{partial}, and related escape hatches in the project sources, and
\path{NonScattering/MainTheorem.lean} runs \texttt{\#print axioms} on the
public theorems.  The pinned-toolchain kernel build has been completed:
\texttt{lake build} succeeds with zero warnings on Lean/mathlib v4.32.1,
and the \texttt{\#print axioms} report shows only the standard classical
axioms (\texttt{propext}, \texttt{Classical.choice},
\texttt{Quot.sound}) for every public theorem.  The build log, the axiom
report, and the pinned transitive dependency manifest are archived in
\path{verification/lean\_kernel\_build\_v1.2.1.log} and
\path{formalization/lake-manifest.json}.

\section{Reproduction}

From the release root, the primary checks are
\begin{lstlisting}[basicstyle=\ttfamily\small,frame=single]
make verify
make lean
make paper
\end{lstlisting}
The verification-only path is
\begin{lstlisting}[basicstyle=\ttfamily\small,frame=single]
cd verification
./reproduce.sh
\end{lstlisting}
and the Lean path is
\begin{lstlisting}[basicstyle=\ttfamily\small,frame=single]
cd formalization
python3 scripts/check_no_holes.py
lake build
\end{lstlisting}
The reproducer authenticates dependencies before executing proof arithmetic.
The root manifest covers every regular release file except generated archives
and the manifest itself.  The release is archived at Zenodo,
\href{https://doi.org/10.5281/zenodo.22016374}{doi:10.5281/zenodo.22016374},
and developed at
\url{https://github.com/Mansehej/non-scattering-schiffer}.

\section{Dependency ledger}

\begin{center}
\small
\begin{tabular}{@{}>{\raggedright\arraybackslash}p{0.30\linewidth}>{\raggedright\arraybackslash}p{0.20\linewidth}>{\raggedright\arraybackslash}p{0.40\linewidth}@{}}
\toprule
Claim & Status & Evidence \\
\midrule
Exact Schiffer zero and infinite tails & Imported theorem & Pinned and
replayed Colbrook--Stepaniants MPFR certificate.\\
Univalence, analytic boundary, non-disc & Imported theorem & Uniform geometry
bounds on the same certified ball.\\
Bessel-deformed equation & Article proof & Direct Laplacian and conformal
calculation.\\
Constants $1906$ and $3$ & Article and supplement proof & Derivation chain
rechecked in exact rationals by Python and Node; terminal scalar
inequalities also in Lean.\\
Uniform fixed point & Article proof & Banach contraction theorem plus the
endpoint certificate; scalar closure rechecked in Python, Node, and
Lean.\\
Scaling and $q\ne1$ & Article proof & Exact dilation and positivity
calculation; scalar part rechecked in Lean.\\
Herglotz and conformal trace identities & Article proof & Standard planar
analysis; outside the Lean scope.\\
Zero scattered field & Article proof & Exterior total field is the incident
field identically; uniqueness of the direct problem cited in the
article.\\
\bottomrule
\end{tabular}
\end{center}

\begin{thebibliography}{CCH16}

\bibitem[Ber80]{Berenstein1980}
Carlos A. Berenstein.
\newblock An inverse spectral theorem and its relation to the {Pompeiu} problem.
\newblock \emph{J. Analyse Math.} 37 (1980), 128--144.

\bibitem[BPS14]{BlastenPaivarintaSylvester2014}
Eemeli Bl{\aa}sten, Lassi P{\"a}iv{\"a}rinta, and John Sylvester.
\newblock Corners always scatter.
\newblock \emph{Comm. Math. Phys.} 331 (2014), no.~2, 725--753.

\bibitem[BST73]{BrownSchreiberTaylor1973}
Leon Brown, Bertram M. Schreiber, and B. Alan Taylor.
\newblock Spectral synthesis and the {Pompeiu} problem.
\newblock \emph{Ann. Inst. Fourier (Grenoble)} 23 (1973), no.~3, 125--154.

\bibitem[CCH16]{CakoniColtonHaddar2016}
Fioralba Cakoni, David Colton, and Houssem Haddar.
\newblock \emph{Inverse Scattering Theory and Transmission Eigenvalues}.
\newblock CBMS-NSF Regional Conference Series in Applied Mathematics 88, SIAM,
Philadelphia, 2016.

\bibitem[CV23]{CakoniVogelius2023}
Fioralba Cakoni and Michael S. Vogelius.
\newblock Singularities almost always scatter: regularity results for
non-scattering inhomogeneities.
\newblock \emph{Comm. Pure Appl. Math.} 76 (2023), no.~12, 4022--4047.

\bibitem[CV26]{CakoniVogelius2026}
Fioralba Cakoni and Michael S. Vogelius.
\newblock Transmission eigenvalues and non-scattering.
\newblock arXiv:2602.06250, 2026.

\bibitem[CVX23]{CakoniVogeliusXiao2023}
Fioralba Cakoni, Michael S. Vogelius, and Jingni Xiao.
\newblock On the regularity of non-scattering anisotropic inhomogeneities.
\newblock \emph{Arch. Ration. Mech. Anal.} (2023), doi:10.1007/s00205-023-01863-y.

\bibitem[CDP26]{CaoLaboraDeDiosPont2026}
Gonzalo Cao-Labora and Jaume de Dios Pont.
\newblock Counterexamples to {Schiffer}'s conjecture.
\newblock arXiv:2608.05114, 2026.

\bibitem[CS26]{ColbrookStepaniants2026}
Matthew J. Colbrook and George Stepaniants.
\newblock A computer-assisted counterexample to the planar {Pompeiu} and
{Schiffer} conjectures.
\newblock arXiv:2608.01579, 2026.

\bibitem[CS26c]{ColbrookStepaniantsCert2026}
Matthew J. Colbrook and George Stepaniants.
\newblock {Pompeiu}--{Schiffer} validation certificate v1.0.0.
\newblock Zenodo, 2026, doi:10.5281/zenodo.21765287.

\bibitem[CK19]{ColtonKress2019}
David Colton and Rainer Kress.
\newblock \emph{Inverse Acoustic and Electromagnetic Scattering Theory}.
\newblock 4th edition, Applied Mathematical Sciences 93, Springer, Cham, 2019.

\bibitem[CM88]{ColtonMonk1988}
David Colton and Peter Monk.
\newblock The inverse scattering problem for time-harmonic acoustic waves in
an inhomogeneous medium.
\newblock \emph{Quart. J. Mech. Appl. Math.} 41 (1988), no.~1, 97--125.

\bibitem[EH15]{ElschnerHu2015}
Johannes Elschner and Guanghui Hu.
\newblock Corners and edges always scatter.
\newblock \emph{Inverse Problems} 31 (2015), no.~1, 015003.

\bibitem[HV25]{HovsepyanVogelius2025}
Narek Hovsepyan and Michael S. Vogelius.
\newblock Scattering from analytic and piecewise analytic inhomogeneities.
\newblock arXiv:2507.13986, 2025.

\bibitem[Kir86]{Kirsch1986}
Andreas Kirsch.
\newblock The denseness of the far field patterns for the transmission problem.
\newblock \emph{IMA J. Appl. Math.} 37 (1986), no.~3, 213--225.

\bibitem[PSV17]{PaivarintaSaloVesalainen2017}
Lassi P{\"a}iv{\"a}rinta, Mikko Salo, and Esa V. Vesalainen.
\newblock Strictly convex corners scatter.
\newblock \emph{Rev. Mat. Iberoam.} 33 (2017), no.~4, 1369--1396.

\bibitem[Pom29]{Pompeiu1929}
Dimitrie Pompeiu.
\newblock Sur certains syst\`emes d'\'equations lin\'eaires et sur une
propri\'et\'e int\'egrale des fonctions de plusieurs variables.
\newblock \emph{C. R. Acad. Sci. Paris} 188 (1929), 1138--1139.

\bibitem[SS21]{SaloShahgholian2021}
Mikko Salo and Henrik Shahgholian.
\newblock Free boundary methods and non-scattering phenomena.
\newblock \emph{Res. Math. Sci.} 8 (2021), no.~4, Paper No. 58.

\bibitem[VX21]{VogeliusXiao2021}
Michael S. Vogelius and Jingni Xiao.
\newblock Finiteness results concerning nonscattering wave numbers for
incident plane and {Herglotz} waves.
\newblock \emph{SIAM J. Math. Anal.} 53 (2021), no.~5, doi:10.1137/20M1367854.

\bibitem[Wil76]{Williams1976}
Stephen A. Williams.
\newblock A partial solution of the {Pompeiu} problem.
\newblock \emph{Math. Ann.} 223 (1976), no.~2, 183--190.

\end{thebibliography}


\end{document}
